\documentclass{article}
\usepackage{graphicx}
\usepackage{amsmath}
\usepackage{amsfonts}
\usepackage{bookmark}
\usepackage{booktabs}
\usepackage{exercises}
\graphicspath{ {./images/} }
\usepackage{indentfirst}
\usepackage{amsthm}
\usepackage{yhmath}
\usepackage{placeins}
\newcommand*{\oldneg}{\mathord{\sim}}
\newcommand\showdiv[1]{\overline{\smash{)}#1}}
\newcommand\Mydiv[1]{%
\smash{\raise.3ex\hbox{$\big)$}}\mkern-8mu
        \overline{\enspace\strut#1}}
\setlength{\tabcolsep}{5pt}
\theoremstyle{definition}
\newtheorem{definition}{Definition}

\setlength{\parindent}{.5in}
\setlength{\parskip}{12pt}
\renewcommand{\arraystretch}{2}
\newcommand{\?}{\stackrel{?}{=}}

\title{4.5 Answers}
\begin{document}
\maketitle

\section*{1}
To condense the expression $\log_3 2x+\log_3 y$, you need to use the proudct property of logarithms.

\section*{2}
One way to evalue $\log_7 12$ using a calculator is by doing exactly that, evaluating $\log_7 12$ on a calculator. Another way is to use the change in base formula 
to determine that \[ \log_7 12=\dfrac{\log 12}{\log 7} \] and then use a calculator for the new expression.

\section*{3}
B

\section*{4}
D

\section*{5}
A

\section*{6}
C

\section*{7}
\[ \log 3=\log{\frac{12}{4}}=\log 12-\log 4=1.079-0.602\approx 0.477 \]

\section*{9}
\[ \log 16=\log (4^2) \] \[ =2\log 4=2(0.602)\approx 1.204 \]

\section*{11}
\[ \log 144=\log (12^2) \] \[ =2\log 12=2(1.079)\approx 2.158 \]

\section*{13}
\[ \log \frac{1}{4}=\log (4^{-1})=-\log 4\approx -0.602 \]

\section*{25}
\[ \ln x^2y^\frac{1}{3}=\ln x^2+\ln y^\frac{1}{3}=2\ln x+\frac{1}{3}\ln y \]

\section*{26}
\[ \log 10x^3=\log 10+\log x^3=1+3\log x \]

\section*{27}
\[ \log_2 \sqrt{x}=\log_2 (x^\frac{1}{2})=\frac{1}{2}\log_2 x \]

\section*{28}
\[ \ln\dfrac{6x^2}{y^4}=\ln 6x^2-\ln y^4=\ln 6+\ln x^2-4\ln y \] \[ =\ln 6+2\ln x-4\ln y \]

\section*{29}
\[ \ln \sqrt[4]{x^3}=\ln (x^\frac{3}{4})=\frac{3}{4}\ln x \]

\section*{30}
\[ \log_3 \sqrt{9x} =\log_3 (3x^\frac{1}{2})=\log_3 3+\log_3 \left( x^\frac{1}{2}\right )=1+\frac{1}{2}\log_3 x \] 

\section*{31}
The error here is that the logarithm product property was used incorrectly. Instead of $\log_a mn=(\log_a m)(\log_a n)$, the property is that 
$\log_a mn=\log_a m+\log_a n$, meaning: \[ \log_2 5x=\log_2 5+\log_2 x \]

\section*{39}
\[ \ln 40+2\ln \frac{1}{2}+\ln x=\ln (40\times\left( \frac{1}{2}\right)^2\times x)=\ln 10x \]

\section*{40}
\[ \log_5 4+\frac{1}{3}\log_5 x=\log_4 4\sqrt[3]{x} \]

\section*{41}
\[ 6\ln 2-4\ln y=\ln 64-\ln y^4 \] \[ =\ln\frac{64}{y^4} \]

\section*{42}
\[ 2(\log_3 20-\log_3 4)+0.5\log_3 4=2\log_3 20-2\log_3 4+\log_3 \left( 4^\frac{1}{2}\right) \]
\[ =\log_3 400-\log_3 16+\log_3 2=\log_3 800-\log_3 16=\log_3\dfrac{800}{16} \]

\section*{43}
C

\section*{44}
D

\section*{56}
\[ \log_8 32=\dfrac{\log_2 32}{\log_2 8}=\frac{5}{3} \]

\section*{58}
\[ \log_2 \frac{15}{7}=\frac{\log\dfrac{15}{7}}{\log{2}}\approx 1.1 \]

\section*{60}
\[ \log_7 \frac{3}{16}=\frac{\log\dfrac{3}{16}}{\log{7}}\approx -0.86 \]

\section*{64}
(a): Let $b=5,N=5,M=4$, meaning: \[ \log_b(M+N)=\log_5(4+5)=\frac{\log 9}{\log 5}\approx 1.37 \] whereas 
\[ \log_b M+\log_b N=\log_5 4+\log_5 5=\frac{\log 4}{\log 5}+1\approx 1.86 \]


(b): The same numbers work, except setting $N=4$ and $M=5$ (because $\log_5(4-5)$ is undefined).

\section*{65}
Let $x=\log_b m$ and $y=\log_b n$. Thus \[ \log_b m+\log_b n=x+y \] Thus $b^x=m$ and $b^y=n$, meaning 
\[ \log_b mn=\log \left( b^x\times b^y\right)=\log_b\left( b^{x+y}\right)=x+y \] thus completing the proof.

\section*{66}
Let $x=\log_b m$ and let $y=\log_b n$. Thus $m=b^x$ and $n=b^y$, meaning \[ \log_b m-\log_b n=x-y \] We then have:
\[ \log_b \frac{m}{n}=\log_b\dfrac{b^x}{b^y}=\log_b\left( b^{x-y}\right)=x-y \] thus completing the proof.

\section*{67}
Let $x=\log_b m$. Thus $m=b^x$ and $m^n=b^{xn}$, meaning:
\[ \log_b \left( m^n\right)=\log_b\left( b^{xn}\right)=xn \] Because $n\log_b m=nx=xn$ by the assumption that $\log_b m=x$, the proof is complete.

\section*{68}
Let $x=\log_c a$, $y=\log_b a$, and $z=\log_b c$. Thus $a=c^x$, $a=b^y$, and $c=b^z$, meaning $c^x=b^y$. We then have:
\[ \frac{\log_b a}{\log_b c}=\frac{\log_b\left( b^y\right)}{\log_b\left( b^z\right)}=\frac{y}{z} \] Next, we know that:
\[ c^x=b^y \] \[ (b^z)^x=b^y \] \[ b^{zx}=b^y \] \[ b^x=b^\frac{y}{z} \] meaning $x=\frac{y}{z}$. 
Because it was assumed that $x=\log_c a$, the proof is complete.

\section*{71}
We know that the function for finding the decibels produced is:
\[ y=10\log\frac{I}{I_0} \] meaning the decibels for the show is:
\[ 10\log\frac{I}{I_0} \] whereas the decibels for the ad is: \[ 10\log\frac{10I}{I_0} \]
To find the difference in decibels, we subract the two quantities:
\[ 10\log\frac{10I}{I_0}-10\log\frac{I}{I_0}=10(\log\frac{10I}{I_0}-\log\frac{I}{I_0}) \]
\[ =10\left( \log 10+\log\frac{I}{I_0}-\log\frac{I}{I_0}\right) \]
\[ =10(1)=10 \]
meaning the tv ad is 10 decibels louder than the show.

\section*{74}
(a): Supose that we are given a wind speed of $2\ln\left( 100h\right)$. If the altitude doubles, then the wind speed would be:
\[ 2\ln\left(100\times 2h\right)=2\ln\left( 2\times 100h\right)=2(\ln 2+\ln\left( 100h\right) ) \]
\[ =2\ln 2+2\ln\left( 100h\right) \]
meaning the wind speed would increase by $2\ln 2=\ln 4$.


(b): \[ 2\ln\left(100h\right)=2\dfrac{\log\left( 100h\right)}{\log e}=\frac{2}{\log e}\left( \log 100+\log h\right) \]
\[ =\frac{2}{\log e}\left( \log h+2\right) \] 

\end{document}
